\section{Descenso de morfismos de preesquemas étales}
\label{IX.3}

\begin{proposition}
\label{IX.3.1}
Sean~$g\colon S'\to S$ un morfismo \emph{sobreyectivo} de preesquemas,
$X$ e~$Y$ dos preesquemas sobre~$S$, y~$X'$, $Y'$ sus preimágenes
sobre~$S'$. Si~$Y$ es no ramificado sobre~$S$, entonces la aplicación
canónica
$$
\Hom_S(X,Y) \to \Hom_{S'}(X',Y')
$$
es inyectiva.
\end{proposition}

En efecto, en virtud de~\Ref{IX.1.6}, un $S$-morfismo
$f\colon X\to Y$ queda determinado cuando se conoce el conjunto subyacente
a su grafo~$\Gamma$, que es una parte de~$Z=X\times_SY$. Como
$$
Z'=Z\times_SS'=X'\times_{S'}Y' \to Z
$$
es sobreyectivo (pues~$S'\to S$ lo es), esta parte~$\Gamma$ queda
determinada cuando se conoce su preimagen en~$X'\times_{S'}Y'$, que no es
otra que el conjunto subyacente al grafo de~$f'$. De ahí la conclusión.

Una parte~$\Gamma$ de~$Z$ es el grafo de un $S$-morfismo
$f\colon X\to Y$ si y solo si es abierta y si el morfismo inducido por
$\pr_1$ de~$\Gamma$ en
\ifthenelse{\boolean{orig}}
{$Y$}
{$X$}
es radicial
\marginpar{234}
y sobreyectivo; \cf \Ref{IX.1}. Cuando se verifica la primera propiedad,
la segunda se verifica si y solo si la preimagen~$\Gamma'$ de~$\Gamma$ en
$Z'$ satisface la misma condición (\Ref{VIII.3.1}). Si, por último, se sabe
que~$Z'\to Z$ es submersivo, lo que sucede en particular si~$S'\to S$ es
universalmente submersivo, entonces~$\Gamma$ es abierta si y solo si lo
es~$\Gamma'$.

Así, el conjunto~$\Hom_S(X,Y)$ está entonces en correspondencia biunívoca
con el conjunto de las partes abiertas~$\Gamma'$ de~$Z'$ tales que el
morfismo de proyección
\ifthenelse{\boolean{orig}}
{$\pr_1\colon Z'\to X'$}
{$\pr_1\colon \Gamma'\to X'$}
sea radicial y sobreyectivo (es decir, que correspondan a un $S'$-morfismo
$f'\colon X'\to Y'$), y que estén saturadas para la relación de
equivalencia definida por~$Z'\to Z$; es decir, cuyas dos preimágenes en
$Z''=Z'\times_Z Z'=Z\times_S S''$ (donde~$S''=S'\times_S S'$), por una y
otra proyección, sean iguales. Ahora bien, estas últimas son los grafos de
los dos $S''$-morfismos~$X''\to Y''$ deducidos de~$f'$ por cambio de base,
mediante una y otra proyección~$S''\to S'$. Hemos obtenido así:

\begin{proposition}
\label{IX.3.2}
Sean~$g\colon S'\to S$ un morfismo \emph{universalmente submersivo} de
preesquemas, $S''=S'\times_S S'$, $X$ e~$Y$ dos $S$-preesquemas, $X'$ e
$Y'$ sus preimágenes sobre~$S'$, y~$X''$, $Y''$ sus preimágenes
sobre~$S''$. Si~$Y$ es étale sobre~$S$, el siguiente diagrama canónico de
aplicaciones es exacto:
$$
\xymatrix@C=.5cm{ {\Hom_S(X,Y)} \ar[r] & {\Hom_{S'}(X',Y')}
\ar@/^2pt/[r]\ar@/_2pt/[r] & {\Hom_{S''}(X'',Y'')}. }
$$
\end{proposition}

Tomando~$X$ e~$Y$ étales sobre~$S$, se obtiene el enunciado siguiente, que,
por lo demás, vuelve a dar~\Ref{IX.3.2} (incluso restringiéndose a~$X=S$,
caso al que siempre puede reducirse~\Ref{IX.3.2} mediante el cambio de base
$X\to S$):

\begin{corollaire}
\label{IX.3.3}
Un morfismo universalmente submersivo de preesquemas es un morfismo de
descenso para la categoría fibrada de los preesquemas étales sobre otros.
\end{corollaire}

Ignoro, por lo demás, si es necesariamente un morfismo de descenso
\emph{efectivo} para la categoría fibrada en cuestión, incluso añadiendo la
hipótesis de que~$S$ sea noetheriano y~$g$ de tipo finito, y limitándose a
los recubrimientos étales. No obstante, en el apartado siguiente daremos
criterios útiles de efectividad.

\begin{corollaire}
\label{IX.3.4}
Sea
\marginpar{235}
$g\colon S'\to S$ un morfismo universalmente submersivo cuyas fibras
$g^{-1}(s)$ sean \og geométricamente conexas\fg,
\index{geometricamente conexo@geométricamente conexo|hyperpage}%
es decir, tales que, para toda extensión~$K/\kres(s)$,
$g^{-1}(s)\otimes_{\kres(s)}K$ sea conexo. Entonces~$S'$ es conexo si~$S$
lo es. El funtor de la categoría de los preesquemas étales sobre~$S$ en la
\ifthenelse{\boolean{orig}}
{categorías}
{categoría}
de los preesquemas étales sobre~$S'$ definido por~$g$ es plenamente fiel.
\end{corollaire}

Una parte de~$S'$ que sea a la vez abierta y cerrada está saturada para la
relación de equivalencia conjuntista definida por~$g$, puesto que las
fibras son conexas; por tanto, es la preimagen de una parte de~$S$, que es
necesariamente abierta y cerrada porque~$g$ es submersivo. Así pues, si~$S$
es conexo, $S'$ lo es. Esto implica también el resultado siguiente: el
compuesto~$fg$ de dos morfismos con fibras universalmente conexas,
\index{universalmente conexo: sinónimo de geométricamente conexo|hyperpage}%
siendo~$f$ universalmente submersivo, tiene fibras universalmente conexas;
si~$S_1'$ y~$S_2'$ sobre~$S$ tienen fibras universalmente conexas, lo mismo
ocurre con~$S_1'\times_S S_2'$. En particular, bajo las condiciones
de~\Ref{IX.3.4}, $S''$ tiene fibras universalmente conexas sobre~$S$. Sean
entonces~$X$ e~$Y$ étales sobre~$S$, y sea~$u'$ un $S'$-morfismo de~$X'$ en
$Y'$; demostremos que es compatible con los datos de descenso (lo cual
entraña la conclusión buscada gracias a~\Ref{IX.3.3}). Sean, en efecto,
$u_1''$ y~$u_2''$ los dos $S''$-morfismos~$X''\to Y''$ deducidos de~$u'$.
El subpreesquema de~$S''$ de coincidencia de~$u_1''$ y~$u_2''$ es un
subpreesquema abierto inducido, cerrado fibra a fibra, como preimagen del
preesquema diagonal de~$Y''$ sobre~$S''$\kern1pt\footnote{Obsérvese que las
fibras de
\ifthenelse{\boolean{orig}}
{$S'$ sobre~$S$}
{$Y$ sobre~$S$}
están separadas.}. Por tanto, es la preimagen de una parte de~$S$. Como
contiene la diagonal en~$S''$, coincide con~$S''$, de donde
$u_1''=u_2''$, cqfd.
